\begin{table}[!h]
\centering
\caption{\label{tab:tipi}Big Five trait scores from the Ten-Item Personality Inventory ($N = 24$).}
\centering
\fontsize{10}{12}\selectfont
\begin{threeparttable}
\begin{tabular}[t]{lcc}
\toprule
Trait & M (SD) & Observed range\\
\midrule
Extraversion & 3.38 (1.60) & 1.0--7.0\\
Agreeableness & 5.83 (1.17) & 3.5--7.0\\
Conscientiousness & 6.02 (1.04) & 4.0--7.0\\
Emotional stability & 4.77 (1.44) & 2.0--7.0\\
Openness & 5.50 (1.48) & 1.5--7.0\\
\bottomrule
\end{tabular}
\begin{tablenotes}
\item \textit{Note.} 
\item Each trait is the mean of two items (one reverse-scored) rated on a 1--7 scale (Gosling et al., 2003). Higher values indicate a stronger expression of the trait. In the \emph{Persuade + Demographics + Personality} condition these scores were passed to the model as part of the participant profile.
\end{tablenotes}
\end{threeparttable}
\end{table}
